% Rédaction : BANDE Hamidou ; relecture : BADO Lionel
\chapter{Supervision avec Prometheus et Grafana}\label{chap:supervision}

\section{Les métriques exposées}
Un middleware Express mesure chaque requête et la bibliothèque \texttt{prom-client} publie les
résultats sur \texttt{/metrics}, au format texte de Prometheus.

\begin{table}[H]
\centering\small
\begin{tabularx}{\linewidth}{@{}llX@{}}
\toprule
\textbf{Métrique} & \textbf{Type} & \textbf{Ce qu'elle mesure} \\
\midrule
\texttt{http\_requests\_total} & Counter & requêtes, par méthode, route et code HTTP \\
\texttt{http\_request\_errors\_total} & Counter & erreurs, classées \texttt{client} (4xx) ou \texttt{serveur} (5xx) \\
\texttt{http\_request\_duration\_seconds} & Histogram & temps de réponse, seaux de 5~ms à 5~s \\
\texttt{http\_requests\_in\_flight} & Gauge & requêtes en cours \\
\texttt{transactions\_operations\_total} & Counter & créations, modifications, suppressions (métrique métier) \\
\bottomrule
\end{tabularx}
\caption{Métriques de l'API : requêtes, erreurs et temps de réponse (exigence E4)}
\end{table}

Un \emph{Counter} ne fait que monter, une \emph{Gauge} monte et descend, un \emph{Histogram} range
les durées par seaux pour calculer des centiles. Le label \texttt{route} contient le \emph{modèle}
de route (\texttt{/api/transactions/:reference}) et jamais l'URL réelle : sinon chaque référence
créerait sa propre série et leur nombre exploserait.

\section{Prometheus}
Prometheus vient lire \texttt{api:3000/metrics} toutes les 5 secondes (modèle \emph{pull}) sous le
nom de tâche \texttt{api-transactions} : l'API n'a pas besoin de connaître Prometheus.

\capture{hamidou_04_prometheus_cibles.png}{Prometheus : la cible \texttt{api-transactions} est UP}{fig:prom-targets}

Trois règles d'alerte sont chargées ; en fonctionnement normal, elles sont toutes inactives.

\begin{table}[H]
\centering\small
\begin{tabularx}{\linewidth}{@{}lXl@{}}
\toprule
\textbf{Alerte} & \textbf{Condition} & \textbf{Durée} \\
\midrule
\texttt{ApiIndisponible} & la cible ne répond plus (\texttt{up == 0}) & 30 s \\
\texttt{TauxErreursServeurEleve} & plus de 5 \% d'erreurs serveur & 1 min \\
\texttt{LatenceP95Elevee} & p95 au-dessus de 500 ms & 2 min \\
\bottomrule
\end{tabularx}
\caption{Règles d'alerte (\texttt{monitoring/prometheus/alert\_rules.yml})}
\end{table}

\capture{hamidou_05_prometheus_alertes.png}{Prometheus : les trois règles d'alerte, inactives}{fig:prom-alertes}

Requêtes PromQL principales :
\begin{lstlisting}
sum(rate(http_requests_total{job="api-transactions"}[1m]))                     # débit
sum(rate(http_request_errors_total[5m])) / sum(rate(http_requests_total[5m]))   # taux d'erreurs
histogram_quantile(0.95, sum by (le) (rate(http_request_duration_seconds_bucket[5m])))  # p95
\end{lstlisting}

\section{Grafana}
Grafana se configure seul au démarrage, à partir des fichiers de \texttt{monitoring/grafana/} :
source de données Prometheus et tableau de bord « API Transactions - Groupe 5 ». Ses 13 panneaux :
API disponible, requêtes par seconde (globales et par route), taux d'erreurs, erreurs par code HTTP,
latence p95, centiles p50/p95/p99, temps de réponse moyen par route, répartition des codes HTTP,
écritures, requêtes en cours, mémoire du processus et latence de la boucle d'événements Node.js.

La figure~\ref{fig:grafana-trafic} a été prise pendant le trafic de démonstration
(\texttt{npm run demo:traffic -- --rps 6}), qui mêle lectures, créations, modifications,
suppressions et erreurs volontaires (400, 404, 409). L'API traite environ 6 requêtes par seconde
avec une latence p95 d'environ 46~ms. Le taux d'erreurs affiché (environ 15~\%) correspond à ces
erreurs \emph{client} volontaires : aucune erreur serveur n'apparaît.

\capture{hamidou_06_grafana_trafic.png}{Le tableau de bord pendant le trafic de démonstration}{fig:grafana-trafic}

\section{Panne contrôlée de la base}
Pendant le trafic, nous avons arrêté PostgreSQL (\texttt{docker compose stop postgres}) pendant
85 secondes, de 17:56:49 à 17:58:14 le 29/10/2026, puis l'avons redémarré
(\texttt{docker compose start postgres}). Observations :
\begin{itemize}
  \item l'API reste vivante : \texttt{/health} répond 200, et \texttt{/health/ready} signale
        la panne (statut \texttt{degrade}, base \texttt{indisponible}) ;
  \item les routes de données répondent 503 \texttt{BASE\_INDISPONIBLE} ; sur la fenêtre de
        6 minutes de la figure~\ref{fig:grafana-panne}, Grafana compte 408 réponses 503 et le taux
        d'erreurs monte à 50~\% ;
  \item les temps de réponse p95 et p99 montent à près de 5 s : c'est le délai d'attente de
        connexion à la base, après lequel l'API répond 503 au lieu de rester bloquée ;
  \item au redémarrage de la base, l'API se reconnecte seule et répond de nouveau 200, sans avoir
        été relancée.
\end{itemize}

\capture{hamidou_07_grafana_panne.png}{Grafana pendant la panne de PostgreSQL : réponses 503 et délais de connexion}{fig:grafana-panne}
